% GENERATED, edit the French unit or the translation table
\mxunit{seconde-s09}{Proportions and percentage change}\mxsubtitle{21--29 January 2027}
\section{Proportion and proportion of a proportion}
\begin{keypoint}[Definitions]
In a nonempty finite set $E$, the proportion of a subset $F$ is $p=\Card(F)/\Card(E)$; $0\leq p\leq1$. Its expression as a percentage is $100p\,\%$.
If $G\subseteq F\subseteq E$, $F\ne\varnothing$, and if $p_1$ is the proportion of $F$ in $E$ and $p_2$ that of $G$ in $F$, the proportion of $G$ in $E$ is $p_1p_2$.
\end{keypoint}
\begin{example}A school has $1500$ students, of whom $652$ are day boarders. Calculate their proportion. In a class, $60\,\%$ of the students play sport; among them, $40\,\%$ are boys. What proportion of the class do the boys who play sport represent?\end{example}
\mxanswerspace[4]{$652/1500\approx0.4347$, that is about $43.47\,\%$. For the class, $0.60\times0.40=0.24$, that is $24\,\%$. The second percentage has the sporty students as its reference, not the whole class.}
\begin{example}In a fictional school of $800$ students, $35\,\%$ play a sport at a club; among them, $45\,\%$ play a team sport. Calculate the two counts and the final proportion.\end{example}
\mxanswerspace[4]{$800\times0.35=280$ students at a club. $280\times0.45=126$ in a team sport. The final proportion is $126/800=0.1575=15.75\,\%$.}
\begin{pitfall}Proportion and rate of change do not answer the same question: a proportion compares a part to a whole, a rate compares a change to an initial value.\end{pitfall}

\mxsection{Absolute change and rate of change}
\begin{keypoint}[Definition]
For an initial value $V_I>0$ and a final value $V_F>0$, the absolute change is $V_F-V_I$; the rate of change is
\[t=\frac{V_F-V_I}{V_I}.\]
The rate has no unit. It is positive for an increase, negative for a decrease. The reference is always the initial value.
\end{keypoint}
\begin{example}A cheese is sold for $2.13$ euros instead of $3.55$ euros. Calculate the absolute change and the rate.\end{example}
\mxanswerspace[3.2]{Absolute change: $2.13-3.55=-1.42$ euro. Rate: $-1.42/3.55=-0.40$, that is a decrease of $40\,\%$.}
\begin{keypoint}[Multiplier]
$CM=V_F/V_I=1+t$, so $V_F=V_I(1+t)$. To recover $V_I$, divide $V_F$ by $1+t$ (nonzero).
\end{keypoint}
\begin{example}A fictional item priced $50000$ Ar increases by $12\,\%$, then separately consider a $12\,\%$ discount on the initial price. Give the final prices. A price after a $20\,\%$ increase is $72000$ Ar; find its initial price.\end{example}
\mxanswerspace[4]{Increase: $50000\times1.12=56000$ Ar. Discount: $50000\times0.88=44000$ Ar. Last question: $V_I=72000/1.20=60000$ Ar.}
\begin{fluency}[5--10 min]Expand $(x-3)^2$; solve $1/x=5$; give the mean of $8;10;12$.\mxanswer{$x^2-6x+9$ ; $x=1/5$ ; $10$.}\end{fluency}

\mxsection{Successive changes}
\begin{keypoint}[Property]
Successive changes with rates $t_1,t_2$ have overall multiplier
\[CM_G=(1+t_1)(1+t_2),\qquad t_G=(1+t_1)(1+t_2)-1.\]
We multiply the multipliers; we do not simply add the rates.
\end{keypoint}
\begin{example}A manufacturer sold $5600$ bicycles in 2022, which is $25\,\%$ more than in 2021. Its sales decrease by $15\,\%$ in 2023. Determine the sales of 2021 and 2023, then the overall rate between these two years.\end{example}
\mxanswerspace[4]{2021: $5600/1.25=4480$ bicycles. 2023: $5600\times0.85=4760$ bicycles. Overall multiplier $1.25\times0.85=1.0625$, so an overall increase of $6.25\,\%$. Check: $(4760-4480)/4480=0.0625$.}
\begin{example}A price increases by $10\,\%$ then decreases by $10\,\%$. Does it return to the initial price?\end{example}
\mxanswerspace[3.2]{$CM_G=1.10\times0.90=0.99$. The final price is $99\,\%$ of the initial price: an overall decrease of $1\,\%$. The two rates refer to different reference values.}
\begin{example}A rate goes from $20\,\%$ to $25\,\%$. Explain the difference between the increase in percentage points and the relative change.\end{example}
\mxanswerspace[3.2]{The gap is $5$ percentage points. The relative change is $(0.25-0.20)/0.20=0.25$, that is an increase of $25\,\%$.}

\mxsection{Reverse change}
\begin{keypoint}[Property]
To return to the initial value after a change of rate $t>-1$, use the multiplier $1/(1+t)$. The reverse rate is
\[t_{\text{reverse}}=\frac1{1+t}-1=\frac{-t}{1+t}.\]
\end{keypoint}
\begin{example}Which rate compensates a $25\,\%$ increase? Which rate compensates a $20\,\%$ decrease?\end{example}
\mxanswerspace[4]{After a $25\,\%$ increase, the reverse multiplier is $1/1.25=0.8$: a $20\,\%$ decrease. After a $20\,\%$ decrease, it is $1/0.8=1.25$: a $25\,\%$ increase.}
\begin{example}A price undergoes a $30\,\%$ discount, then a $10\,\%$ increase. What additional percentage increase restores the initial price?\end{example}
\mxanswerspace[5]{The current multiplier is $0.70\times1.10=0.77$. The required multiplier is $1/0.77=100/77$. The rate is $23/77\approx0.2987$, that is about $29.87\,\%$. Check $0.77\times(100/77)=1$.}
\begin{teachernote}Have students state the base of the percentage at each step. The reverse-rate formula is justified by a product equal to $1$; memorisation alone is fragile.\end{teachernote}
